\documentclass{article}
\usepackage[utf8]{inputenc}
\usepackage{amsmath}
\usepackage{amssymb}
\usepackage{geometry}
\geometry{letterpaper, margin=1in}
\usepackage{hyperref}

\title{Homework 1.1: Introduction to Differential Equations - Problem Sets}
\author{}
\date{}

\begin{document}

\maketitle
\tableofcontents
\newpage

\section{Concept Checklist \& Question Mapping}
This section contains all the concepts and problem types extracted from the source material. The generated problems ensure comprehensive coverage of all these specific ideas.

\begin{itemize}
    \item \textbf{Concept 1: Equation Anatomy and Classification (Variables \& ODE vs PDE)} \\
    \textit{Skill:} Identifying independent and dependent variables, and distinguishing Ordinary Differential Equations (ODEs) from Partial Differential Equations (PDEs). \\
    \textit{Mapped Questions:} 1, 2, 3, 4, 5, 6, 7, 8, 9, 10
    
    \item \textbf{Concept 2: Determining Order} \\
    \textit{Skill:} Identifying the highest derivative present in the equation. \\
    \textit{Mapped Questions:} 1, 2, 3, 4, 5, 6, 7, 8, 9, 10
    
    \item \textbf{Concept 3: Determining Linearity} \\
    \textit{Skill:} Checking if the dependent variable and its derivatives appear only to the first power, are not multiplied together, and do not appear inside non-linear functions (like sine, square roots, etc.). \\
    \textit{Mapped Questions:} 1, 2, 3, 4, 5, 6, 7, 8, 9, 10
    
    \item \textbf{Concept 4: Mathematical Translation \& Modeling} \\
    \textit{Skill:} Translating English descriptions of physical situations (rates of change, proportionality, inverse proportionality) into differential equations. \\
    \textit{Mapped Questions:} 11, 12, 13, 14, 15, 16
    
    \item \textbf{Concept 5: Kinematic Root-Finding (Physics Application)} \\
    \textit{Skill:} Setting up systems of equations for constant acceleration ($d = v_0t + \frac{1}{2}at^2$), taking advantage of algebraic division to eliminate unknowns, and using square roots of ratios to solve for time. \\
    \textit{Mapped Questions:} 17, 18, 19, 20
    
    \item \textbf{Concept 6: Verifying Solutions to Initial Value Problems (IVPs)} \\
    \textit{Skill:} Differentiating a given function and substituting it back into a differential equation to prove it forms a true mathematical statement. \\
    \textit{Mapped Questions:} 21, 22, 23, 24, 25
    
    \item \textbf{Concept 7: Advanced Diagnostic Testing ("Find the Flaw")} \\
    \textit{Skill:} Identifying subtle conceptual errors in proposed solutions regarding order vs. degree, identifying independent vs. dependent variables, and correctly applying mathematical terminology. \\
    \textit{Mapped Questions:} 26, 27, 28, 29, 30
\end{itemize}

\newpage

\section{Problem Sets}

\subsection*{Part A: Equation Anatomy and Classification}
For problems 1--10, classify the given differential equation by identifying: 
(a) the independent variable(s), (b) the dependent variable(s), (c) whether it is an ODE or PDE, (d) the order of the equation, and (e) whether it is linear or nonlinear.

\begin{enumerate}
    \item $5 \frac{d^3 y}{dx^3} - 2x \frac{dy}{dx} + 7y = \ln(x)$
    \item $\frac{du}{dt} = u \cos(t) + u^2$
    \item $\frac{\partial^2 T}{\partial x^2} + \frac{\partial^2 T}{\partial y^2} = \frac{\partial T}{\partial t}$
    \item $y'' + \sin(y) = 0$
    \item $\left( \frac{dy}{dx} \right)^2 + y = x$
    \item $\frac{dx}{dt} + \frac{dy}{dt} = 5t$
    \item $t \frac{d^2 x}{dt^2} = \frac{1}{x}$
    \item $e^x \frac{d^4 y}{dx^4} = y$
    \item $y y''' - x y = 0$
    \item $\sqrt{1 - x^2} \frac{dy}{dx} - y = 0$
\end{enumerate}

\subsection*{Part B: Mathematical Translation and Modeling}
For problems 11--16, write a differential equation that fits the physical description. Use $k$ for the constant of proportionality unless otherwise specified.

\begin{enumerate}
    \setcounter{enumi}{10}
    \item The rate of change of the area $A$ of a healing wound with respect to time $t$ is proportional to the square root of the area $A$.
    \item The rate of change of volume $V$ of a sphere with respect to its radius $r$ is proportional to the square of its radius.
    \item Newton's Law of Cooling states that the rate of change of the temperature $T$ of an object with respect to time $t$ is proportional to the difference between $T$ and the ambient room temperature $T_r$.
    \item The rate of change of a commodity's price $P$ with respect to time $t$ is inversely proportional to the price $P$.
    \item The acceleration of a particle (the rate of change of its velocity $v$ with respect to time $t$) is proportional to the cube of its velocity.
    \item The rate of growth of a bacteria population $B$ over time $t$ is proportional to the product of $B$ and the remaining carrying capacity, given by $(1000 - B)$.
\end{enumerate}

\subsection*{Part C: Kinematic Root-Finding}
For problems 17--20, solve the kinematics scenarios. Both competitors begin from a standing start ($v_0 = 0$) and proceed with a constant acceleration.

\begin{enumerate}
    \setcounter{enumi}{16}
    \item Racer 1 and Racer 2 are participating in a drag race. Racer 1 covers the last $\frac{1}{5}$ of the track in $4$ seconds. Racer 2 covers the last $\frac{1}{3}$ of the track in $6$ seconds. Who wins the race and by how much time? (Round to the nearest thousandth of a second).
    \item Car A and Car B are on a test track. Car A covers the last $\frac{1}{4}$ of the distance in $2$ seconds. Car B covers the last $\frac{1}{2}$ of the distance in $5$ seconds. Who wins the race and by how much time? (Round to the nearest thousandth of a second).
    \item Runner X and Runner Y are competing in a sprint. Runner X covers the last $\frac{1}{9}$ of the distance in $1$ second. Runner Y covers the last $\frac{1}{10}$ of the distance in $1.5$ seconds. Who wins the race and by how much time? (Round to the nearest thousandth of a second).
    \item Cyclist M and Cyclist N race on a velodrome. Cyclist M covers the last $\frac{1}{6}$ of the distance in $3$ seconds. Cyclist N covers the last $\frac{1}{8}$ of the distance in $2$ seconds. Who wins the race and by how much time? (Round to the nearest thousandth of a second).
\end{enumerate}

\subsection*{Part D: Verifying Solutions}
For problems 21--25, verify that the given function is a solution to the corresponding differential equation. Show all steps of differentiation and substitution.

\begin{enumerate}
    \setcounter{enumi}{20}
    \item Verify that $y = 3e^{-4x}$ is a solution to the ODE: $y' + 4y = 0$.
    \item Verify that $y = \cos(3t)$ is a solution to the ODE: $y'' + 9y = 0$.
    \item Verify that $y = \frac{1}{x}$ is a solution to the ODE: $x y' + y = 0$.
    \item Verify that $y = x^3$ is a solution to the ODE: $x^2 y'' - 6y = 0$.
    \item Verify that $N = 10e^{2t}$ is a solution to the ODE: $\frac{dN}{dt} - 2N = 0$.
\end{enumerate}

\subsection*{Part E: Advanced Diagnostic Testing ("Find the Flaw")}
For problems 26--30, you are presented with a problem and a flawed solution. Your task is to identify the conceptual flaw and provide the correct answer.

\begin{enumerate}
    \setcounter{enumi}{25}
    \item \textbf{Problem:} Classify the equation $x^3 y'' + y = 0$ for linearity. \\
    \textbf{Flawed Solution:} Because the variable $x$ is raised to the 3rd power, the equation contains a nonlinear term. Therefore, the equation is nonlinear.
    
    \item \textbf{Problem:} Find the order of the differential equation $\frac{d^2 y}{dx^2} + \left(\frac{dy}{dx}\right)^5 = x$. \\
    \textbf{Flawed Solution:} We examine the derivatives. The first derivative term is raised to the 5th power. Because 5 is the highest power of any derivative present, the order of this differential equation is 5.
    
    \item \textbf{Problem:} Write a differential equation for: "The rate of change of mass $M$ with respect to time $t$ is inversely proportional to time $t$." \\
    \textbf{Flawed Solution:} Rate of change of $M$ is $\frac{dM}{dt}$. Inversely proportional to $t$ means we subtract $t$ from a constant. Thus, the equation is $\frac{dM}{dt} = k - t$.
    
    \item \textbf{Problem:} Classify the equation $\frac{\partial^2 u}{\partial x^2} + \frac{\partial^2 u}{\partial y^2} = 0$ as an ODE or a PDE. \\
    \textbf{Flawed Solution:} The equation contains two dependent variables, $x$ and $y$. Because there is more than one dependent variable, this must be a Partial Differential Equation (PDE).
    
    \item \textbf{Problem:} Classify the equation $y \frac{dy}{dx} = x$ for linearity. \\
    \textbf{Flawed Solution:} We look at the dependent variable $y$. It appears to the first power. Its derivative $\frac{dy}{dx}$ also appears to the first power. Because there are no powers greater than 1 on $y$ or its derivatives, the equation is linear.
\end{enumerate}

\newpage

\section{Detailed Solutions}

\subsection*{Part A Solutions}
\begin{enumerate}
    \item $5 \frac{d^3 y}{dx^3} - 2x \frac{dy}{dx} + 7y = \ln(x)$ \\
    \textbf{Variables:} Independent: $x$, Dependent: $y$. \\
    \textbf{ODE/PDE:} Ordinary derivatives only. \textbf{ODE}. \\
    \textbf{Order:} Highest derivative is the third derivative. \textbf{Order 3}. \\
    \textbf{Linearity:} $y$ and its derivatives are to the first power, not multiplied together. (Note: $\ln(x)$ is a function of the independent variable, which is perfectly fine). \textbf{Linear}.
    
    \item $\frac{du}{dt} = u \cos(t) + u^2$ \\
    \textbf{Variables:} Independent: $t$, Dependent: $u$. \\
    \textbf{ODE/PDE:} \textbf{ODE}. \\
    \textbf{Order:} \textbf{Order 1}. \\
    \textbf{Linearity:} The term $u^2$ violates the rule that the dependent variable must be to the first power. \textbf{Nonlinear}.
    
    \item $\frac{\partial^2 T}{\partial x^2} + \frac{\partial^2 T}{\partial y^2} = \frac{\partial T}{\partial t}$ \\
    \textbf{Variables:} Independent: $x, y, t$. Dependent: $T$. \\
    \textbf{ODE/PDE:} Partial derivatives ($\partial$) are used. \textbf{PDE}. \\
    \textbf{Order:} Highest derivative is second order. \textbf{Order 2}. \\
    \textbf{Linearity:} $T$ and all its partial derivatives are to the first power. \textbf{Linear}.
    
    \item $y'' + \sin(y) = 0$ \\
    \textbf{Variables:} Independent: $x$ (implied), Dependent: $y$. \\
    \textbf{ODE/PDE:} \textbf{ODE}. \\
    \textbf{Order:} \textbf{Order 2}. \\
    \textbf{Linearity:} The dependent variable $y$ is inside a transcendental function ($\sin$). \textbf{Nonlinear}.
    
    \item $\left( \frac{dy}{dx} \right)^2 + y = x$ \\
    \textbf{Variables:} Independent: $x$, Dependent: $y$. \\
    \textbf{ODE/PDE:} \textbf{ODE}. \\
    \textbf{Order:} \textbf{Order 1}. \\
    \textbf{Linearity:} The first derivative is squared. \textbf{Nonlinear}.
    
    \item $\frac{dx}{dt} + \frac{dy}{dt} = 5t$ \\
    \textbf{Variables:} Independent: $t$, Dependent: $x, y$. \\
    \textbf{ODE/PDE:} There is only one independent variable ($t$). Therefore, this is an \textbf{ODE} (specifically a system of ODEs). \\
    \textbf{Order:} \textbf{Order 1}. \\
    \textbf{Linearity:} $x$ and $y$ and their derivatives are to the first power. \textbf{Linear}.
    
    \item $t \frac{d^2 x}{dt^2} = \frac{1}{x}$ \\
    \textbf{Variables:} Independent: $t$, Dependent: $x$. \\
    \textbf{ODE/PDE:} \textbf{ODE}. \\
    \textbf{Order:} \textbf{Order 2}. \\
    \textbf{Linearity:} The dependent variable is in the denominator ($x^{-1}$), which violates the first power rule. \textbf{Nonlinear}.
    
    \item $e^x \frac{d^4 y}{dx^4} = y$ \\
    \textbf{Variables:} Independent: $x$, Dependent: $y$. \\
    \textbf{ODE/PDE:} \textbf{ODE}. \\
    \textbf{Order:} \textbf{Order 4}. \\
    \textbf{Linearity:} The term $e^x$ is a coefficient containing only the independent variable. $y$ and its derivative are to the first power. \textbf{Linear}.
    
    \item $y y''' - x y = 0$ \\
    \textbf{Variables:} Independent: $x$, Dependent: $y$. \\
    \textbf{ODE/PDE:} \textbf{ODE}. \\
    \textbf{Order:} \textbf{Order 3}. \\
    \textbf{Linearity:} The dependent variable $y$ is multiplied by its derivative $y'''$. \textbf{Nonlinear}.
    
    \item $\sqrt{1 - x^2} \frac{dy}{dx} - y = 0$ \\
    \textbf{Variables:} Independent: $x$, Dependent: $y$. \\
    \textbf{ODE/PDE:} \textbf{ODE}. \\
    \textbf{Order:} \textbf{Order 1}. \\
    \textbf{Linearity:} The coefficient $\sqrt{1-x^2}$ depends only on $x$. The variable $y$ and its derivative are to the first power. \textbf{Linear}.
\end{enumerate}

\subsection*{Part B Solutions}
\begin{enumerate}
    \setcounter{enumi}{10}
    \item "Rate of change of $A$" translates to $\frac{dA}{dt}$. "Proportional to" means $= k \cdot (...)$. "Square root of $A$" means $\sqrt{A}$. \\
    \textbf{Equation:} $\frac{dA}{dt} = k\sqrt{A}$.
    
    \item "Rate of change of $V$ with respect to $r$" translates to $\frac{dV}{dr}$. \\
    \textbf{Equation:} $\frac{dV}{dr} = k r^2$.
    
    \item "Rate of change of $T$" translates to $\frac{dT}{dt}$. "Difference between $T$ and $T_r$" translates to $(T - T_r)$. \\
    \textbf{Equation:} $\frac{dT}{dt} = k(T - T_r)$.
    
    \item "Rate of change of $P$" translates to $\frac{dP}{dt}$. "Inversely proportional to $P$" means dividing a constant by $P$. \\
    \textbf{Equation:} $\frac{dP}{dt} = \frac{k}{P}$.
    
    \item Acceleration is $\frac{dv}{dt}$. "Proportional to the cube of its velocity" translates to $k v^3$. \\
    \textbf{Equation:} $\frac{dv}{dt} = k v^3$.
    
    \item "Rate of growth of $B$" translates to $\frac{dB}{dt}$. Product means multiplication. \\
    \textbf{Equation:} $\frac{dB}{dt} = k B(1000 - B)$.
\end{enumerate}

\subsection*{Part C Solutions}
\textit{General setup for kinematics (Problems 17-20):} Let $D$ be total distance. $d(t) = \frac{1}{2}at^2$. If a racer covers the last fraction $f$ of the distance in $T$ seconds, they cover the first $(1-f)$ of the distance in $(t - T)$ seconds. 
Setting up the ratio: $\frac{(1-f)D}{D} = \frac{\frac{1}{2}a(t-T)^2}{\frac{1}{2}at^2} \implies 1-f = \left(\frac{t-T}{t}\right)^2$. 
Taking the square root: $\sqrt{1-f} = 1 - \frac{T}{t} \implies t = \frac{T}{1 - \sqrt{1-f}}$.

\begin{enumerate}
    \setcounter{enumi}{16}
    \item \textbf{Racer 1:} Last $\frac{1}{5}$ of track in $4$ seconds. $f = \frac{1}{5}$, $1-f = \frac{4}{5}$, $T=4$.
    $$ t_1 = \frac{4}{1 - \sqrt{4/5}} = \frac{4}{1 - 2/\sqrt{5}} = \frac{4\sqrt{5}}{\sqrt{5} - 2} = 4\sqrt{5}(\sqrt{5}+2) = 20 + 8\sqrt{5} \approx 37.889 \text{ s} $$
    \textbf{Racer 2:} Last $\frac{1}{3}$ of track in $6$ seconds. $f = \frac{1}{3}$, $1-f = \frac{2}{3}$, $T=6$.
    $$ t_2 = \frac{6}{1 - \sqrt{2/3}} \approx \frac{6}{1 - 0.81650} \approx 32.702 \text{ s} $$
    Because $32.702 < 37.889$, \textbf{Racer 2 wins}. The difference is $37.889 - 32.702 = 5.187$ seconds.
    
    \item \textbf{Car A:} Last $\frac{1}{4}$ of track in $2$ seconds. $f = \frac{1}{4}$, $1-f = \frac{3}{4}$, $T=2$.
    $$ t_A = \frac{2}{1 - \sqrt{3/4}} = \frac{2}{1 - \sqrt{3}/2} = \frac{4}{2 - \sqrt{3}} = 4(2+\sqrt{3}) \approx 14.928 \text{ s} $$
    \textbf{Car B:} Last $\frac{1}{2}$ of track in $5$ seconds. $f = \frac{1}{2}$, $1-f = \frac{1}{2}$, $T=5$.
    $$ t_B = \frac{5}{1 - \sqrt{1/2}} = \frac{5\sqrt{2}}{\sqrt{2} - 1} = 5\sqrt{2}(\sqrt{2}+1) = 10 + 5\sqrt{2} \approx 17.071 \text{ s} $$
    Because $14.928 < 17.071$, \textbf{Car A wins}. The difference is $17.071 - 14.928 = 2.143$ seconds.
    
    \item \textbf{Runner X:} Last $\frac{1}{9}$ of track in $1$ second. $f = \frac{1}{9}$, $1-f = \frac{8}{9}$, $T=1$.
    $$ t_X = \frac{1}{1 - \sqrt{8/9}} = \frac{3}{3 - \sqrt{8}} = 3(3+\sqrt{8}) = 9 + 6\sqrt{2} \approx 17.485 \text{ s} $$
    \textbf{Runner Y:} Last $\frac{1}{10}$ of track in $1.5$ seconds. $f = \frac{1}{10}$, $1-f = \frac{9}{10}$, $T=1.5$.
    $$ t_Y = \frac{1.5}{1 - \sqrt{9/10}} = \frac{1.5\sqrt{10}}{\sqrt{10} - 3} = 1.5\sqrt{10}(\sqrt{10}+3) = 15 + 4.5\sqrt{10} \approx 29.230 \text{ s} $$
    Because $17.485 < 29.230$, \textbf{Runner X wins}. The difference is $29.230 - 17.485 = 11.745$ seconds.
    
    \item \textbf{Cyclist M:} Last $\frac{1}{6}$ of track in $3$ seconds. $f = \frac{1}{6}$, $1-f = \frac{5}{6}$, $T=3$.
    $$ t_M = \frac{3}{1 - \sqrt{5/6}} = \frac{3\sqrt{6}}{\sqrt{6} - \sqrt{5}} = 3\sqrt{6}(\sqrt{6}+\sqrt{5}) = 18 + 3\sqrt{30} \approx 34.431 \text{ s} $$
    \textbf{Cyclist N:} Last $\frac{1}{8}$ of track in $2$ seconds. $f = \frac{1}{8}$, $1-f = \frac{7}{8}$, $T=2$.
    $$ t_N = \frac{2}{1 - \sqrt{7/8}} = \frac{2}{1 - \frac{\sqrt{14}}{4}} = \frac{8}{4 - \sqrt{14}} = \frac{8(4+\sqrt{14})}{16-14} = 4(4+\sqrt{14}) = 16 + 4\sqrt{14} \approx 30.967 \text{ s} $$
    Because $30.967 < 34.431$, \textbf{Cyclist N wins}. The difference is $34.431 - 30.967 = 3.464$ seconds.
\end{enumerate}

\subsection*{Part D Solutions}
\begin{enumerate}
    \setcounter{enumi}{20}
    \item \textbf{Verify:} $y = 3e^{-4x}$ into $y' + 4y = 0$. \\
    Derivative: Using chain rule, $y' = 3(-4)e^{-4x} = -12e^{-4x}$. \\
    Substitute: $(-12e^{-4x}) + 4(3e^{-4x}) = -12e^{-4x} + 12e^{-4x} = 0$. \\
    $0 = 0$, the solution is verified.
    
    \item \textbf{Verify:} $y = \cos(3t)$ into $y'' + 9y = 0$. \\
    Derivatives: $y' = -3\sin(3t)$, $y'' = -9\cos(3t)$. \\
    Substitute: $(-9\cos(3t)) + 9(\cos(3t)) = 0$. \\
    $0 = 0$, the solution is verified.
    
    \item \textbf{Verify:} $y = \frac{1}{x} = x^{-1}$ into $x y' + y = 0$. \\
    Derivative: $y' = -1x^{-2} = -\frac{1}{x^2}$. \\
    Substitute: $x(-\frac{1}{x^2}) + \frac{1}{x} = -\frac{1}{x} + \frac{1}{x} = 0$. \\
    $0 = 0$, the solution is verified.
    
    \item \textbf{Verify:} $y = x^3$ into $x^2 y'' - 6y = 0$. \\
    Derivatives: $y' = 3x^2$, $y'' = 6x$. \\
    Substitute: $x^2(6x) - 6(x^3) = 6x^3 - 6x^3 = 0$. \\
    $0 = 0$, the solution is verified.
    
    \item \textbf{Verify:} $N = 10e^{2t}$ into $\frac{dN}{dt} - 2N = 0$. \\
    Derivative: $\frac{dN}{dt} = 20e^{2t}$. \\
    Substitute: $(20e^{2t}) - 2(10e^{2t}) = 20e^{2t} - 20e^{2t} = 0$. \\
    $0 = 0$, the solution is verified.
\end{enumerate}

\subsection*{Part E Solutions}
\begin{enumerate}
    \setcounter{enumi}{25}
    \item \textbf{Identify the Flaw:} The solution incorrectly applies linearity tests to the independent variable $x$. Linearity only depends on the dependent variable $y$. \\
    \textbf{Correction:} The variables $y$ and $y''$ are strictly to the first power and are not multiplied together. The $x^3$ acts as a valid coefficient. Thus, the equation is \textbf{linear}.
    
    \item \textbf{Identify the Flaw:} The flaw is confusing the algebraic degree (the exponent 5) with the order of the derivative. \\
    \textbf{Correction:} Order is determined by the highest derivative taken. The highest derivative is $\frac{d^2 y}{dx^2}$, so the \textbf{order is 2}. (The exponent 5 merely makes it nonlinear).
    
    \item \textbf{Identify the Flaw:} The flaw represents a misunderstanding of mathematical vocabulary. "Inversely proportional" does not mean subtraction. \\
    \textbf{Correction:} Inversely proportional means proportional to the reciprocal (division). The correct equation is $\frac{dM}{dt} = \frac{k}{t}$.
    
    \item \textbf{Identify the Flaw:} The flaw confuses independent and dependent variables. $x$ and $y$ are in the denominator of the derivative operator, meaning they are the independent variables. $u$ is the single dependent variable. \\
    \textbf{Correction:} An equation is a PDE because there are \textbf{multiple independent variables} ($x$ and $y$). 
    
    \item \textbf{Identify the Flaw:} The flaw is ignoring the product rule of linearity; the dependent variable cannot be multiplied by itself or its derivatives. \\
    \textbf{Correction:} The term $y \frac{dy}{dx}$ is a product of the dependent variable and its derivative. This explicitly violates linearity. Thus, the equation is \textbf{nonlinear}.
\end{enumerate}

\newpage

\section{Concept Checklist (Review)}
\textit{Use this list to review specific concepts by referencing their corresponding problem numbers.}

\begin{itemize}
    \item \textbf{Concept 1: Equation Anatomy and Classification (Variables \& ODE vs PDE)} $\rightarrow$ See Q: 1-10
    \item \textbf{Concept 2: Determining Order} $\rightarrow$ See Q: 1-10
    \item \textbf{Concept 3: Determining Linearity} $\rightarrow$ See Q: 1-10
    \item \textbf{Concept 4: Mathematical Translation \& Modeling} $\rightarrow$ See Q: 11-16
    \item \textbf{Concept 5: Kinematic Root-Finding (Physics Application)} $\rightarrow$ See Q: 17-20
    \item \textbf{Concept 6: Verifying Solutions to Initial Value Problems (IVPs)} $\rightarrow$ See Q: 21-25
    \item \textbf{Concept 7: Advanced Diagnostic Testing ("Find the Flaw")} $\rightarrow$ See Q: 26-30
\end{itemize}

\end{document}